\section{From design to implementation}
\label{part:implementation}

Every preceding part describes something that does not exist. This part is about
the one piece that does, what it measured, and why the measurement matters more
than the six parts before it.

\subsection{What actually exists}

It is worth being exact, because the distinction between a specification and a
system is the easiest one to lose while writing a document like this.

\begin{center}
\begin{tabular}{@{}l r p{0.44\linewidth}@{}}
\toprule
Artifact & Size & Kind \\
\midrule
MLIR fixtures            & 47   & test data \\
C provenance files       & 42   & test data \\
Actor examples           & 7    & test data \\
Integration tests        & 8    & test data \\
\midrule
\code{check\_harness.py} & 929  & metadata \emph{validation}; checks fixtures are
                                  well-formed. Not analysis \\
\code{refusal\_rate.py}  & 179  & a scanner over IR text. Not analysis \\
\code{sps-scan.cpp}      & 174  & \textbf{an actual analysis}: an MLIR
                                  \code{SparseDataFlowAnalysis} with a
                                  \textsf{Low}/\textsf{High} lattice \\
\bottomrule
\end{tabular}
\end{center}

So: roughly a hundred test files, and of thirteen hundred lines of code exactly
one hundred and seventy-four implement an information-flow analysis. Everything
else validates, measures, or documents.

\begin{remark}[Why the accounting matters]
It is easy to write ``the architecture refuses 63\% of functions'' when what was
measured is ``63\% of functions contain an operation with no source location.''
The first is a claim about a system; the second is a claim about clang's output
plus a rule from a design document. Only the second was measured. Keeping the
two apart is the difference between a result and a rationalisation.
\end{remark}

\subsection{The first real measurement}

\code{sps-scan} was run over every fixture whose policy survives parsing --- 17
of 47, the rest encoding their policy in SSA names that \code{mlir-opt}
discards. It emits findings with a stable reason and a source location:

\begin{center}
\begin{minipage}{0.9\linewidth}\small
\begin{verbatim}
[secret-dependent-branch] llvm.cond_br @ loc(...:67:5)
[secret-to-public-sink]   llvm.store   @ loc(...:69:5)
findings: 2
\end{verbatim}
\end{minipage}
\end{center}

The results sort into four groups, and all four are informative.

\paragraph{Caught, correctly.} \fx{predecessor\_choice\_blockarg.bad} yields two
findings, including the MT-CM6 case where a secret selects a merge block argument
with no secret on any SSA operand edge. \fx{prefix\_causal\_release.bad} yields
one. These are the countermodels a re-implementation would most plausibly fail,
and a 174-line pass catches them.

\paragraph{Silent, correctly.} Five fixtures yield zero findings: the
public-allocation control, the offset-disjoint control, the overwritten-slot
control, the repaired error oracle, and the constant-latency wolfSSL helper
profile. A pass that reported everything as unsafe --- the failure mode the
corpus was previously unable to detect at all --- fails here.

\paragraph{Imprecise, as predicted.} Two precision controls fire, one finding
each: the identical-successor branch and the value-cancelling \code{xor}.

This is \textbf{not} a defect. Those fixtures exist to document that a forward
\textsf{Low}/\textsf{High} lattice \emph{must} lose precision at exactly those
sites, and each carries \code{l1 disposition: relational-required} saying so. The
implementation is behaving the way \cref{step:diag} says it is required to. It is
the first time a prediction in this document was confirmed by running something.

\paragraph{Missed, for known reasons.} Three bad fixtures yield zero: the
audience mismatch, the error oracle, and the affected wolfSSL helper profile. They need audience reasoning, release-policy conformance and a
helper-latency contract respectively --- none of which a taint pass implements.
These are the shape of the remaining work, and they are visible only because
something ran.

\subsection{The argument for building before designing further}

This document has been wrong, and corrected by measurement, six times:

\begin{center}
\begin{tabular}{@{}p{0.43\linewidth} p{0.45\linewidth}@{}}
\toprule
Claim & Corrected by \\
\midrule
whole-run release equality as the
    product's initial constraint
  & the spec forbids it; the encoding
    is unsound in the permissive
    direction \\
the declassification marker as an
    opaque function
  & attribute inference deducing
    \code{returned} once the definition
    is co-visible \\
per-site quantification in the
    decision rule
  & a build reporting \Verified{} under
    it and refusing under universal
    discharge \\
a register-only reduction of the laundering case
  & it did not reproduce; the conversion needs a memory operand \\
an arithmetic-mask repair
  & folded straight back into a select, byte-identical to the leaky twin \\
a refusal surface of about 1\%
  & 4.8\% per observation, and 63\% per function on real code \\
\bottomrule
\end{tabular}
\end{center}

Six for six. Not one of these was caught by review, by internal consistency, or
by re-reading the specification; every one was caught by running a tool and
looking at the output. The corrected versions are better, but the pattern is the
point: \textbf{in this problem domain, a design claim that has not been executed
is not evidence.}

The remaining untested claims are all design claims. So the highest-value next
action is not more design --- it is to make the smallest real thing falsifiable.

\subsection{The smallest useful next step}

\code{sps-scan} is not referenced anywhere in the harness. Wiring it in converts
the whole corpus from a specification into a test of an implementation:

\begin{enumerate}[nosep,leftmargin=1.6em]
\item Add it to the diagnostic \code{RUN} of the satisfiable bring-up gates. That
      turns ``the analysis must emit this reason here'' from an aspiration into a
      pass/fail.
\item Add it to the four precision controls under their declared
      \RelReq{} disposition. This is the first false-positive regression net the
      corpus has ever had --- and the run above shows exactly which two controls
      it will legitimately trip.
\item Re-derive the refusal number from the analysis rather than from a text
      scanner, with a real policy scoping the observation set.
\end{enumerate}

None of this needs the lowering architecture of Part~\ref{part:lowering}, the
\code{.ll}-versus-MLIR decision, or the thirty unmigrated fixtures. Those gate
later layers. The diagnostic layer is buildable now, against a corpus that is
already known to exercise it.

\begin{remark}[What is durable]
The architecture in this document has been rewritten twice in a day and will
likely change again. What has not changed, and will survive any rewrite, is the
set of reproduced compiler behaviours, the fixtures that encode them, and the
measurement scripts. Those are the assets. A design that cannot yet be executed
should be weighted accordingly --- including this one.
\end{remark}
